%%% Local Variables:
%%% mode: latex
%%% TeX-master: t
%%% End:

% ============================================================
% 第五章：参考文献
% 本章介绍 BibTeX 文献管理和引用方法。
% ============================================================

\chapter{参考文献}\label{chap:references}

\section{参考文献格式规范}

本院格式规范要求：
\begin{compactitem}
    \item 参考文献按在正文中出现的顺序编号，如 \cite{MARTIN20131497}。
    \item 引用标号置于所引内容最后一个字右上角。
    \item 所有被引用文献必须列入参考文献，同一篇文献只用一个序号。
    \item 参考文献内容使用小四号宋体，1.25倍行距。
\end{compactitem}

\section{BibTeX 文件格式}

参考文献存放在 \texttt{bib/refs.bib} 文件中，每条文献是一个条目，格式如下：

\subsection{期刊论文（article）}

\begin{verbatim}
@article{引用key,
    author  = {作者1 and 作者2 and 作者3},
    title   = {论文题目},
    journal = {期刊名},
    year    = {2023},
    volume  = {100},
    number  = {1},
    pages   = {1--10},
    doi     = {10.xxxx/xxxxx}
}
\end{verbatim}

示例（本模板 \texttt{bib/refs.bib} 中已有的条目）：
\begin{verbatim}
@article{MARTIN20131497,
    author  = {M.J. Martin},
    title   = {Nuclear Data Sheets for A = 152},
    journal = {Nuclear Data Sheets},
    year    = {2013},
    volume  = {114},
    pages   = {1497--1847}
}
\end{verbatim}

\subsection{学位论文（phdthesis / mastersthesis）}

\begin{verbatim}
@phdthesis{张三2023,
    author = {张三},
    title  = {基于XXXX方法的核反应截面研究},
    school = {中国原子能科学研究院},
    year   = {2023},
    type   = {博士学位论文}
}

@mastersthesis{李四2022,
    author = {李四},
    title  = {XXX实验方案设计},
    school = {中国原子能科学研究院},
    year   = {2022}
}
\end{verbatim}

\subsection{书籍（book）}

\begin{verbatim}
@book{Blatt1952,
    author    = {Blatt, John M. and Weisskopf, Victor F.},
    title     = {Theoretical Nuclear Physics},
    publisher = {Springer},
    address   = {New York},
    year      = {1952}
}
\end{verbatim}

\subsection{会议论文（inproceedings）}

\begin{verbatim}
@inproceedings{Wang2022,
    author    = {Wang, X. and Zhang, Y.},
    title     = {Title of the Paper},
    booktitle = {Proceedings of the International Conference on ...},
    year      = {2022},
    pages     = {100--105}
}
\end{verbatim}

\section{引用命令}

本模板提供两种引用命令，对应院格式规范中的两种标注方式：

\begin{compactitem}
    \item \verb|\cite{key}|：行内引用，编号以方括号形式嵌入正文，
          效果如 \cite{KIBEDI200235}，适合"由文献[X]可知……"的写法。
    \item \verb|\upcite{key}|：上标引用，编号以上标方括号形式出现，
          效果如\upcite{MARTIN20131497}，适合在句末或词后标注，
          是本院格式规范中最常用的引用形式。
\end{compactitem}

\verb|\upcite| 是本模板在 \texttt{main.tex} 中自定义的命令：
\begin{verbatim}
\newcommand{\upcite}[1]{\textsuperscript{\textsuperscript{\cite{#1}}}}
\end{verbatim}

引用多篇文献时，用英文逗号分隔多个 key，编号将自动合并：
\begin{verbatim}
\cite{key1,key2,key3}      % 行内，效果：[1,2,3]
\upcite{key1,key2,key3}    % 上标，效果：[1,2,3]（上标）
\end{verbatim}
实际效果：\cite{MARTIN20131497,KIBEDI200235}（行内）；\upcite{MARTIN20131497,KIBEDI200235}（上标）。

格式规范示例：

\begin{verbatim}
% 上标引用（推荐，句末标注）
二次铣削\upcite{ref1}可以显著提高表面质量。
原位生成的TiB主要有针状或晶须状\upcite{ref2,ref3}。

% 行内引用（用于"由文献[X]可知"句式）
由文献\cite{ref4,ref5}可知，蠕变断裂以沿晶断裂为主。
\end{verbatim}

\section{如何获取 BibTeX 条目}

不需要手动输入每条文献，可以从数据库直接导出：

\subsection{从 Google Scholar 导出}

\begin{compactenum}
    \item 搜索文献，点击文献下方的引号图标（引用）。
    \item 在弹出窗口底部点击"BibTeX"。
    \item 复制全部内容，粘贴到 \texttt{bib/refs.bib} 文件中。
\end{compactenum}

\subsection{从 CNKI（知网）导出}

\begin{compactenum}
    \item 搜索文献，勾选需要的文献。
    \item 点击"导出与分析" → "导出文献"。
    \item 选择"BibTeX"格式，下载并合并到 \texttt{refs.bib}。
\end{compactenum}

\subsection{从 Web of Science / Scopus 导出}

\begin{compactenum}
    \item 选中文献，点击"导出"（Export）。
    \item 选择"BibTeX"格式并下载。
\end{compactenum}

\subsection{使用 JabRef 管理文献}

JabRef 是一款开源的 BibTeX 文献管理工具，可图形化管理 \texttt{.bib} 文件，
支持从 DOI、arXiv、ISBN 自动填充文献信息，强烈推荐使用。

\section{引用 key 的命名建议}

BibTeX 的 key（如 \texttt{MARTIN20131497}）用于在正文中引用，
建议采用统一的命名规则：

\begin{compactitem}
    \item 英文文献：\texttt{作者姓+年份}，如 \texttt{Blatt1952}
    \item 中文文献：\texttt{作者姓+年份}，如 \texttt{张三2023}
    \item 同一作者同一年多篇：加字母区分，如 \texttt{Wang2022a}、\texttt{Wang2022b}
\end{compactitem}

\textbf{注意}：key 不能含有中文，不能含有空格，区分大小写。

\section{参考文献样式}

本模板默认使用 GB/T 7714 顺序编码制（数字引用）：

\begin{verbatim}
% main.tex 中的设置（默认）
\bibliographystyle{gbt7714-numerical}
\bibliography{bib/refs}
\end{verbatim}

如需切换为著者-出版年制（如"[张三, 2023]"形式），修改为：

\begin{verbatim}
\bibliographystyle{ustcauthoryear}
\end{verbatim}

\section{本章小结}

本章介绍了 BibTeX 文件的格式、各类文献的写法、引用命令的使用，
以及从主流数据库导出文献条目的方法。
第~\ref{chap:advanced}~章将介绍代码块、算法和其他高级排版功能。
